% =============================================================================
% §6 Numerical Study  --  budget ~11 pp (less whatever moved to §3's numerical
% verification section)
% rubric: Project Body, results & discussion (/20--30)
%
% This chapter now reads purely as the descriptor method's own numerical study;
% the published collocation route's reproduction and characterisation has its
% own numerical verification section in §3.
%
% Self-containment rule: every number quoted here must be derivable from what is
% printed in the thesis or its appendices.  Catalogue of all cases -> Appendix G.
% =============================================================================

\subsection{Experimental Design and Cost Conventions}
% State BEFORE the first comparison: what counts as a gate, the assumed
% decomposition, how ancilla are counted, and which direction the convention
% biases the comparison.  Also the two scope statements:
%   (i) all results are classical simulations, dim <~ 2000 capped by simulation
%       cost, never by the algorithm;
%   (ii) all solution errors are classical state-vector reads, no measurement
%       protocol is executed.
% Software (DESolverLib / desolver_hpc / Setonix) named here, detail -> Appendix I.

\subsection{Encoding Cost of the Descriptor Construction}
% Headline table: gates, ancilla, alpha, exactness, gap -- DESCRIPTOR SIDE ONLY.
% The published/collocation baseline for the same quantities is characterised in
% §3's numerical verification section; cross-reference it here rather than
% repeating that table -- present only the descriptor numbers plus the delta
% against the §3 baseline.

\subsection{Spectral Gap and Success Probability}

\subsection{Solution Accuracy on Representative Problems}
% THREE cases carried in depth, no more.  Everything else -> Appendix G table.

\subsection{Nonlinear Benchmarks}

\subsection{An Application Case Study}
% ONE of Black--Scholes or Navier--Stokes.  Not both.

\subsection{Summary of Findings}
% 4--6 numbered findings, at least one a limitation of the descriptor construction
% itself (Hilbert-space growth; the doubled-space route projecting, not preparing).
